\documentclass[book.tex]{subfiles}

\begin{document}
\chapter{Orbital Angular Momentum}

\label{chap:orbital_angular_momentum}
\noindent\rule{11cm}{0.4pt} \\
\textbf{Prerequisites:} \autoref{chap:spin} \\
\textbf{Difficulty Level:} ** \\
\noindent\rule{11cm}{0.4pt}
\vspace{5mm}

In this chapter, we define about the orbital angular momentum operator $L$ and total angular momentum operator $J$. Recall that classical angular momentum is a vector $\vec{L} \in \mathbb{R}^3$ is defined by the cross product equation

\begin{align}
    \vec{L} &= \vec{r} \times \vec{p}
\end{align}

Where $\vec{r} \in \mathbb{R}^3$ is the position vector and $\vec{p} \in \mathbb{R}^3$ is the momentum vector. To define the analogous quantum mechanical operation, we convert these quantities into operators.

\begin{align}
    L &= r \times p
\end{align}

Here $r$ and $p$ are the position and momentum operators. Recall that these are vector operators

\begin{align}
    r &= (r_x, r_y, r_z) \\
    p &= (p_x, p_y, p_z)
\end{align}
Where $r_x$, $p_x$ etc are operators. Component wise, the vector is given by
\begin{align}
    L &= (L_x, L_y, L_z) \\
    &= (r_y p_z - r_z p_y, r_x p_z - r_z p_x, r_x p_y - r_y p_x)
\end{align}
Note that $L_x$, $L_y$, and $L_z$ are themselves operators.

\section{Commutation Relations}

The angular momentum operators satisfy the following commutation relationships

\begin{align}
    [L_x, L_y] &= i \hbar L_z \\
    [L_y, L_z] &= i \hbar L_x \\
    [L_z, L_x] &= i \hbar L_y
\end{align}

We will prove the first of these commutation relationships explicitly and leave the other relationships to the exercises.

\begin{align}
    [L_x, L_y] &= [r_y p_z - r_z p_y, ]
\end{align}

\section{Exercises}
\begin{enumerate}
    \item Prove the other two commutation relationships
    \begin{align}
    [L_y, L_z] &= i \hbar L_x \\
    [L_z, L_x] &= i \hbar L_y
    \end{align}
\end{enumerate}

\end{document}